% Zdroj: PDF priklady-leto-2023.pdf, fyzická str. 7-8. Značka: otázka studijního zaměření I3. Zařazeno: I3.
\section{Komponentový systém}
\szzotazka{léto 2023}{18}{Komponentový systém (herní framework, generika, komponenty)}

\noindent\textit{(Pozn.: v~rámci současné akreditace nešlo o~společnou otázku (q1--q4), ale o~otázku jiné specializace (specializační sekce, ne UI). Do společného okruhu I3 tedy formálně nepatří; tematicky sem ovšem spadá a~je výborná k~procvičení.)}\par

\noindent\textbf{Zadání.}\par
Naprogramujte v C++, C\# nebo Javě všechny níže zmíněné a další potřebné typy (včetně
zmíněných metod) pro implementaci následujícího komponentového herního frameworku.
Ve hře vytvořené pomocí vašeho frameworku má každý objekt ve scéně (např. ležící kámen,
nebo odpočívající trol) být reprezentovaný instancí typu \texttt{GameObject} (dále jen
GO). Samotný GO má reprezentovat množinu komponent -- kde komponenta je pro nás
libovolný vhodný typ (dle vašeho návrhu), jehož instance přidává k instanci GO nějaká
silně typovaná data, a/nebo nějakou další funkcionalitu. Od každého druhu komponenty
může být na instanci GO připojena právě 0 nebo 1 instance daného druhu komponenty.
Každá instance nějaké komponenty může být připojena na maximálně 1 GO.

\begin{enumerate}
\item Na GO by měla být vhodná metoda pro připojení instance nějaké komponenty ke GO --
  pokud GO komponentu stejného druhu již má, tak nová instance nahradí instanci původní.

  Na GO by měla být vhodná metoda pro vrácení instance komponenty zvoleného druhu
  (požadovaný druh si nějak vhodně může zvolit volající). Pokud na GO žádná instance
  požadovaného typu komponenty připojená není, tak to má funkce nějak rozumně indikovat
  (funkci budeme chtít běžně používat pro detekci, zda daný typ komponenty na GO je,
  nebo není).

  Na GO má být metoda \texttt{Update} s parametrem \texttt{timeDelta} jako 32-bitové
  float číslo, které reprezentuje čas v sekundách od předchozího volání metody
  \texttt{Update} na stejném GO. Metoda \texttt{Update} na GO má projít všechny instance
  připojených komponent, a pro komponenty, které mají metodu \texttt{Update} se stejným
  parametrem, se má tato zavolat. Pořadí volání metod \texttt{Update} na komponentách GO
  není důležité. Ve svém objektovém návrhu vhodně zohledněte, že autor komponenty má mít
  možnost se rozhodnout, zda komponenta bude metodu \texttt{Update} mít nebo nikoliv.

\item Napište implementaci komponenty \texttt{Position}, která ke GO přidává 32-bitové
  float souřadnice X a Y reprezentující polohu GO na 2D herní ploše. Napište implementaci
  komponenty \texttt{Health}, která ke GO přidává 32-bitové float hodnoty maximálního
  a aktuálního zdraví (aktuální zdraví je vždy v rozsahu 0 až max včetně).

  Napište příklad kódu, který vytvoří instanci GO pro kámen, který má komponentu
  \texttt{Position}, a instanci GO pro trola, který má komponenty \texttt{Position}
  a \texttt{Health} (max $=30$, aktuální $=20$).

\item Doplňte vhodné komponenty reprezentující efekt otravy a efekt krvácení -- oba
  efekty vyžadují, aby GO, ke kterému jsou připojené, měl komponentu \texttt{Health}
  (pokud GO takovou komponentu nemá, tak efekt nic nedělá). Dále jsou oba efekty
  parametrizované 32-bit float hodnotou, která reprezentuje, kolik aktuálního zdraví
  efekt odečítá každou sekundu. Pokud je efekt připojený k nějakému GO, tak má dojít
  k úpravě zdraví GO dle parametru efektu. Na jednom GO může být maximálně 1 efekt
  otravy a nezávisle maximálně 1 efekt krvácení (tedy GO může být bez efektu, nebo mít
  1 otravu a žádné krvácení, nebo žádnou otravu a 1 krvácení, nebo 1 otravu a 1 krvácení).
  Do budoucna budeme přidávat další podobné efekty.

  Napište implementaci funkce, která dostane seznam GO a hodnotu zranění za sekundu --
  funkce projde všechny předané GO a pokud GO má komponentu \texttt{Health}, tak ke GO
  připojí novou instanci efektu otravy nebo krvácení parametrizované předanou hodnotou
  zranění za sekundu. Volající musí mít navíc možnost si nějak vhodně zvolit, zda funkce
  aplikuje na předané GO efekt otravy nebo efekt krvácení.

  Napište příklad kódu, který na výše vytvořený kámen a trola zkusí aplikovat efekt
  krvácení s hodnotou zranění 1.2 za sekundu.
\end{enumerate}

\medskip
\noindent\textbf{Náčrt řešení.}\par
\textit{(V~PDF bez náčrtu řešení; náčrt doplněn k~výkladu okruhu: vlastní vzorové řešení v~C++, ověřené překladem a~během.)}

\medskip
Klíčová myšlenka: \texttt{GameObject} drží komponenty v~mapě klíčované \emph{typem} komponenty (\texttt{std::\allowbreak type\_index}). Tím je „nejvýš jedna instance daného druhu`` zdarma a vrácení komponenty podle typu je triviální. Volitelnou metodu \texttt{Update} řeší zvláštní rozhraní \texttt{IUpdatable}.

\textbf{(1) Komponenta, rozhraní \texttt{IUpdatable} a \texttt{GameObject}:}
\begin{lstlisting}[language=C++]
struct Component  { virtual ~Component() = default; };
struct IUpdatable {                                     // volitelne chovani
    virtual ~IUpdatable() = default;
    virtual void Update(float dt) = 0;
};

class GameObject {
    std::unordered_map<std::type_index, std::unique_ptr<Component>> comps_;
public:
    template <typename T> void Add(std::unique_ptr<T> c) {
        comps_[std::type_index(typeid(T))] = std::move(c);  // nahradi stary druh
    }
    template <typename T> T* Get() const {
        auto it = comps_.find(std::type_index(typeid(T)));
        return it == comps_.end() ? nullptr               // indikace "nemam"
                                  : static_cast<T*>(it->second.get());
    }
    void Update(float dt) {
        for (auto& [type, c] : comps_)
            if (auto* u = dynamic_cast<IUpdatable*>(c.get()))  // jen kdo Update ma
                u->Update(dt);
    }
};
\end{lstlisting}

\textbf{(2) Konkrétní komponenty a scéna (kámen, trol):}
\begin{lstlisting}[language=C++]
struct Position : Component { float x = 0, y = 0; };
class Health : public Component {             // invariant: 0 <= cur_ <= max_
    float max_, cur_;
public:
    Health(float max, float cur) : max_(max), cur_(std::clamp(cur, 0.0f, max)) {}
    float Max() const { return max_; }
    float Current() const { return cur_; }
    void SetCurrent(float v) { cur_ = std::clamp(v, 0.0f, max_); }  // hlida rozsah
};

auto kamen = std::make_unique<GameObject>();
kamen->Add(std::make_unique<Position>());
auto trol = std::make_unique<GameObject>();
trol->Add(std::make_unique<Position>());
trol->Add(std::make_unique<Health>(30.0f, 20.0f));   // max = 30, aktualni = 20
\end{lstlisting}

\textbf{(3) Efekty.} Efekt je komponenta \emph{a zároveň} \texttt{IUpdatable}; v~\texttt{Update} si najde \texttt{Health} svého GO a ubírá zdraví. Otrava a krvácení sdílejí logiku, ale musejí to být \emph{dva různé typy} (mapa klíčuje podle typu, jinak by se přepsaly), aby mohly být na GO obě naráz. Volbu efektu předáme \emph{šablonovým} parametrem typu:
\begin{lstlisting}[language=C++]
class Effect : public Component, public IUpdatable {
protected:
    GameObject* owner_;  float perSec_;     // pozorovatel vlastnika + parametr
public:
    Effect(GameObject* o, float p) : owner_(o), perSec_(p) {}
    void Update(float dt) override {
        if (auto* hp = owner_->Get<Health>())            // nic, chybi-li Health
            hp->SetCurrent(hp->Current() - perSec_ * dt);
    }
};
struct Poison   : Effect { using Effect::Effect; };  // dva ruzne typy,
struct Bleeding : Effect { using Effect::Effect; };  // aby sly na GO obe naraz

template <typename E>                          // volba efektu = typovy parametr
void applyEffect(const std::vector<GameObject*>& gos, float perSec) {
    for (GameObject* go : gos)
        if (go->Get<Health>())                 // jen na GO s komponentou Health
            go->Add(std::make_unique<E>(go, perSec));
}
// pouziti: na kamen a trola zkusime krvaceni 1.2/s (kamen nema Health -> preskoci)
applyEffect<Bleeding>({kamen.get(), trol.get()}, 1.2f);
\end{lstlisting}

\noindent Klíčování podle typu (\texttt{type\_index}) splní „nejvýš jedna instance daného druhu`` zadarmo; držení přes \texttt{unique\_ptr} zaručí, že komponenta patří nejvýš jednomu GO; zapouzdřená \texttt{Health} hlídá invariant $0 \le$ aktuální $\le$ max ze zadání na jediném místě; \texttt{Get<T>} vrací \texttt{nullptr} při absenci (vhodná indikace pro dotaz „má GO tuto komponentu?``); volitelnost \texttt{Update} přes \texttt{IUpdatable} a \texttt{dynamic\_cast} respektuje, že ne každá komponenta \texttt{Update} má. Volba efektu v~šablonovém parametru je statický polymorfismus; alternativou je běhový příznak nebo tovární funkce.

\medskip
\begin{csalt}[komponenty přes generika]
Komponentový systém je idiomatický v~C\# (styl Unity). \texttt{Get<T>()} přes generika
s~omezením vrací komponentu nebo \texttt{null}; volba efektu je generický typový parametr.
Bez hlaviček, bez instanciačních chyb šablon, bez ruční správy paměti:
\begin{lstlisting}[style=cs]
class GameObject {
    readonly Dictionary<Type, Component> comps = new();
    public T? Get<T>() where T : Component               // null, chybi-li komponenta
        => comps.TryGetValue(typeof(T), out var c) ? (T)c : null;
    public void Add(Component c) => comps[c.GetType()] = c;   // klic = skutecny typ
}
abstract class Effect : Component, IUpdatable {
    public GameObject Owner { get; set; } = null!;
    public float PerSec { get; set; }
    public void Update(float dt) { var h = Owner.Get<Health>(); if (h != null) h.Current -= PerSec * dt; }
}
sealed class Poison : Effect { }  sealed class Bleeding : Effect { }
void ApplyEffect<E>(IEnumerable<GameObject> gos, float perSec) where E : Effect, new() {
    foreach (var go in gos)
        if (go.Get<Health>() != null)                   // jen na GO s komponentou Health
            go.Add(new E { Owner = go, PerSec = perSec });
}
\end{lstlisting}
Klíčování slovníkem podle \texttt{Type} dá „nejvýš jedna instance druhu`` zdarma; volitelný
\texttt{Update} řeší rozhraní \texttt{IUpdatable} a~test \texttt{if (c is IUpdatable u)} místo \texttt{dynamic\_cast}.
\texttt{Health.Current} je vlastnost, jejíž \texttt{set} ořízne hodnotu do rozsahu 0 až max (není vypsáno).
\end{csalt}
